%%%%%%%%%%%%%%%%%%%%%%%%%%%%%%%%%%%%%%%%%%%%%%%%%%%%%%%%%%%%
% 13.4 CONTINUUM MECHANICS
% The Composition of Advanced Mathematics
%%%%%%%%%%%%%%%%%%%%%%%%%%%%%%%%%%%%%%%%%%%%%%%%%%%%%%%%%%%%

\section{Continuum Mechanics}
\label{sec:continuum-mechanics}

\prerequisite{
Multivariable calculus, vector calculus, linear algebra, differential
equations, partial differential equations, tensor notation, mathematical
physics, and mathematical modeling.
}

\learningobjectives{
By the end of this section, the reader should be able to:
\begin{itemize}
    \item explain the continuum hypothesis;
    \item distinguish material and spatial descriptions;
    \item define reference and current configurations;
    \item define the motion map of a continuum;
    \item understand deformation;
    \item construct the deformation gradient;
    \item interpret the Jacobian determinant;
    \item understand local changes in length, area, and volume;
    \item distinguish displacement, velocity, and acceleration fields;
    \item understand material derivatives;
    \item distinguish Lagrangian and Eulerian descriptions;
    \item derive conservation of mass;
    \item understand the continuity equation;
    \item formulate linear momentum balance;
    \item understand body forces and surface forces;
    \item define traction vectors;
    \item understand Cauchy's stress theorem;
    \item define the Cauchy stress tensor;
    \item interpret normal and shear stresses;
    \item understand stress symmetry;
    \item formulate angular momentum balance;
    \item understand strain measures;
    \item define infinitesimal strain;
    \item understand finite deformation;
    \item recognize the right Cauchy--Green tensor;
    \item recognize the left Cauchy--Green tensor;
    \item define the Green--Lagrange strain tensor;
    \item understand principal strains and principal stresses;
    \item recognize eigenvalue problems in continuum mechanics;
    \item understand constitutive laws;
    \item distinguish kinematic, balance, and constitutive equations;
    \item understand linear elasticity;
    \item formulate Hooke's law in three dimensions;
    \item understand isotropic elasticity;
    \item recognize Lamé parameters;
    \item understand Young's modulus and Poisson ratio;
    \item derive the Navier--Cauchy equations;
    \item understand incompressibility;
    \item understand volumetric and deviatoric deformation;
    \item recognize hydrostatic and deviatoric stress;
    \item understand strain energy;
    \item recognize hyperelasticity conceptually;
    \item understand viscoelasticity conceptually;
    \item recognize Maxwell and Kelvin--Voigt models;
    \item understand plasticity conceptually;
    \item recognize yield criteria;
    \item understand thermal strain;
    \item formulate energy balance;
    \item understand constitutive restrictions conceptually;
    \item recognize the second law of thermodynamics in continuum models;
    \item understand boundary and initial conditions;
    \item recognize weak and variational formulations;
    \item understand finite element ideas conceptually;
    \item recognize wave propagation in elastic media;
    \item understand longitudinal and transverse waves;
    \item understand dimensional analysis in continuum mechanics;
    \item recognize nondimensional parameters;
    \item understand model assumptions and limitations;
    \item connect continuum mechanics with fluid dynamics, elasticity,
          thermodynamics, PDEs, numerical methods, and mathematical physics.
\end{itemize}
}

%%%%%%%%%%%%%%%%%%%%%%%%%%%%%%%%%%%%%%%%%%%%%%%%%%%%%%%%%%%%
\subsection{What Is Continuum Mechanics?}
\label{subsec:what-is-continuum-mechanics}
%%%%%%%%%%%%%%%%%%%%%%%%%%%%%%%%%%%%%%%%%%%%%%%%%%%%%%%%%%%%

Continuum mechanics studies materials by treating matter as continuously
distributed through space.

Instead of tracking individual atoms or molecules, the model assigns fields
such as

\[
\boxed{
\rho(\mathbf{x},t)
}
\]

for density,

\[
\boxed{
\mathbf{v}(\mathbf{x},t)
}
\]

for velocity, and

\[
\boxed{
\boldsymbol{\sigma}(\mathbf{x},t)
}
\]

for stress.

This approximation is effective when the physical scale of interest is much
larger than the microscopic molecular scale.

%%%%%%%%%%%%%%%%%%%%%%%%%%%%%%%%%%%%%%%%%%%%%%%%%%%%%%%%%%%%
\subsection{Continuum Hypothesis}
\label{subsec:continuum-hypothesis}
%%%%%%%%%%%%%%%%%%%%%%%%%%%%%%%%%%%%%%%%%%%%%%%%%%%%%%%%%%%%

The continuum hypothesis assumes that physical quantities vary continuously
enough to be represented by mathematical fields.

For example, density is treated as

\[
\boxed{
\rho
=
\rho(\mathbf{x},t)
}
\]

rather than as a collection of discrete particles.

The continuum approximation allows the use of differential equations.

%%%%%%%%%%%%%%%%%%%%%%%%%%%%%%%%%%%%%%%%%%%%%%%%%%%%%%%%%%%%
\subsection{Material Particles}
\label{subsec:material-particles}
%%%%%%%%%%%%%%%%%%%%%%%%%%%%%%%%%%%%%%%%%%%%%%%%%%%%%%%%%%%%

A material particle is an idealized infinitesimal element of matter.

It is identified by its position in a reference configuration:

\[
\boxed{
\mathbf{X}.
}
\]

As the continuum moves, that same material particle occupies a current
position

\[
\boxed{
\mathbf{x}.
}
\]

%%%%%%%%%%%%%%%%%%%%%%%%%%%%%%%%%%%%%%%%%%%%%%%%%%%%%%%%%%%%
\subsection{Reference and Current Configurations}
\label{subsec:reference-current-configurations}
%%%%%%%%%%%%%%%%%%%%%%%%%%%%%%%%%%%%%%%%%%%%%%%%%%%%%%%%%%%%

The reference configuration describes the body at a chosen initial or
undeformed state.

Write

\[
\boxed{
\mathbf{X}
\in
\Omega_0.
}
\]

The current configuration at time \(t\) is

\[
\boxed{
\mathbf{x}
\in
\Omega_t.
}
\]

The mapping between them is the motion.

%%%%%%%%%%%%%%%%%%%%%%%%%%%%%%%%%%%%%%%%%%%%%%%%%%%%%%%%%%%%
\subsection{Motion Map}
\label{subsec:continuum-motion-map}
%%%%%%%%%%%%%%%%%%%%%%%%%%%%%%%%%%%%%%%%%%%%%%%%%%%%%%%%%%%%

The motion of the body is represented by

\[
\boxed{
\mathbf{x}
=
\boldsymbol{\chi}(\mathbf{X},t).
}
\]

For each material point \(\mathbf{X}\), the function
\(\boldsymbol{\chi}\) gives its location at time \(t\).

Thus

\[
\boxed{
\boldsymbol{\chi}
:
\Omega_0
\to
\Omega_t.
}
\]

%%%%%%%%%%%%%%%%%%%%%%%%%%%%%%%%%%%%%%%%%%%%%%%%%%%%%%%%%%%%
\subsection{Displacement}
\label{subsec:continuum-displacement}
%%%%%%%%%%%%%%%%%%%%%%%%%%%%%%%%%%%%%%%%%%%%%%%%%%%%%%%%%%%%

The displacement field is

\[
\boxed{
\mathbf{u}(\mathbf{X},t)
=
\boldsymbol{\chi}(\mathbf{X},t)
-
\mathbf{X}.
}
\]

Therefore,

\[
\boxed{
\mathbf{x}
=
\mathbf{X}
+
\mathbf{u}(\mathbf{X},t).
}
\]

Displacement measures the change in material position relative to the
reference configuration.

%%%%%%%%%%%%%%%%%%%%%%%%%%%%%%%%%%%%%%%%%%%%%%%%%%%%%%%%%%%%
\subsection{Velocity}
\label{subsec:continuum-velocity}
%%%%%%%%%%%%%%%%%%%%%%%%%%%%%%%%%%%%%%%%%%%%%%%%%%%%%%%%%%%%

The material velocity is

\[
\boxed{
\mathbf{v}
=
\frac{\partial\boldsymbol{\chi}}
{\partial t}
\bigg|_{\mathbf{X}}.
}
\]

In the Lagrangian description,

\[
\mathbf{v}
=
\mathbf{v}(\mathbf{X},t).
\]

In the Eulerian description,

\[
\boxed{
\mathbf{v}
=
\mathbf{v}(\mathbf{x},t).
}
\]

%%%%%%%%%%%%%%%%%%%%%%%%%%%%%%%%%%%%%%%%%%%%%%%%%%%%%%%%%%%%
\subsection{Acceleration}
\label{subsec:continuum-acceleration}
%%%%%%%%%%%%%%%%%%%%%%%%%%%%%%%%%%%%%%%%%%%%%%%%%%%%%%%%%%%%

The acceleration of a material particle is

\[
\boxed{
\mathbf{a}
=
\frac{D\mathbf{v}}{Dt}.
}
\]

In Eulerian coordinates,

\[
\boxed{
\frac{D\mathbf{v}}{Dt}
=
\frac{\partial\mathbf{v}}{\partial t}
+
(\mathbf{v}\cdot\nabla)\mathbf{v}.
}
\]

The second term is the convective acceleration.

%%%%%%%%%%%%%%%%%%%%%%%%%%%%%%%%%%%%%%%%%%%%%%%%%%%%%%%%%%%%
\subsection{Material Derivative}
\label{subsec:material-derivative}
%%%%%%%%%%%%%%%%%%%%%%%%%%%%%%%%%%%%%%%%%%%%%%%%%%%%%%%%%%%%

For a scalar field

\[
f(\mathbf{x},t),
\]

the material derivative is

\[
\boxed{
\frac{Df}{Dt}
=
\frac{\partial f}{\partial t}
+
\mathbf{v}\cdot\nabla f.
}
\]

This gives the rate of change experienced by a moving material particle.

%%%%%%%%%%%%%%%%%%%%%%%%%%%%%%%%%%%%%%%%%%%%%%%%%%%%%%%%%%%%
\subsection{Worked Example: Material Derivative}
\label{subsec:worked-material-derivative}
%%%%%%%%%%%%%%%%%%%%%%%%%%%%%%%%%%%%%%%%%%%%%%%%%%%%%%%%%%%%

\begin{workedexamplebox}{Temperature Experienced by a Moving Particle}

Suppose

\[
T(x,t)=x^2+t
\]

and velocity is

\[
v(x,t)=2.
\]

Then

\[
\frac{\partial T}{\partial t}
=
1,
\]

and

\[
\frac{\partial T}{\partial x}
=
2x.
\]

Therefore,

\[
\frac{DT}{Dt}
=
1
+
2(2x).
\]

Hence

\begin{answerbox}
\[
\boxed{
\frac{DT}{Dt}
=
1+4x.
}
\]
\end{answerbox}

\end{workedexamplebox}

%%%%%%%%%%%%%%%%%%%%%%%%%%%%%%%%%%%%%%%%%%%%%%%%%%%%%%%%%%%%
\subsection{Lagrangian Description}
\label{subsec:lagrangian-description-continuum}
%%%%%%%%%%%%%%%%%%%%%%%%%%%%%%%%%%%%%%%%%%%%%%%%%%%%%%%%%%%%

The Lagrangian or material description follows individual material points.

Independent variables are

\[
\boxed{
\mathbf{X},t.
}
\]

Fields are written, for example,

\[
\boxed{
\mathbf{x}
=
\boldsymbol{\chi}(\mathbf{X},t).
}
\]

This viewpoint is especially natural in solid mechanics.

%%%%%%%%%%%%%%%%%%%%%%%%%%%%%%%%%%%%%%%%%%%%%%%%%%%%%%%%%%%%
\subsection{Eulerian Description}
\label{subsec:eulerian-description-continuum}
%%%%%%%%%%%%%%%%%%%%%%%%%%%%%%%%%%%%%%%%%%%%%%%%%%%%%%%%%%%%

The Eulerian or spatial description observes what occurs at fixed spatial
locations.

Independent variables are

\[
\boxed{
\mathbf{x},t.
}
\]

Fields such as

\[
\boxed{
\rho(\mathbf{x},t)
}
\]

and

\[
\boxed{
\mathbf{v}(\mathbf{x},t)
}
\]

describe material currently passing through location \(\mathbf{x}\).

This viewpoint is especially natural in fluid mechanics.

%%%%%%%%%%%%%%%%%%%%%%%%%%%%%%%%%%%%%%%%%%%%%%%%%%%%%%%%%%%%
\subsection{Deformation Gradient}
\label{subsec:deformation-gradient}
%%%%%%%%%%%%%%%%%%%%%%%%%%%%%%%%%%%%%%%%%%%%%%%%%%%%%%%%%%%%

The deformation gradient is

\[
\boxed{
\mathbf{F}
=
\frac{\partial\mathbf{x}}
{\partial\mathbf{X}}.
}
\]

In components,

\[
\boxed{
F_{iJ}
=
\frac{\partial x_i}
{\partial X_J}.
}
\]

Using

\[
\mathbf{x}
=
\mathbf{X}+\mathbf{u},
\]

we obtain

\[
\boxed{
\mathbf{F}
=
\mathbf{I}
+
\nabla_{\mathbf{X}}\mathbf{u}.
}
\]

%%%%%%%%%%%%%%%%%%%%%%%%%%%%%%%%%%%%%%%%%%%%%%%%%%%%%%%%%%%%
\subsection{Local Line Deformation}
\label{subsec:local-line-deformation}
%%%%%%%%%%%%%%%%%%%%%%%%%%%%%%%%%%%%%%%%%%%%%%%%%%%%%%%%%%%%

An infinitesimal reference line element

\[
d\mathbf{X}
\]

transforms into

\[
\boxed{
d\mathbf{x}
=
\mathbf{F}\,d\mathbf{X}.
}
\]

Thus \(\mathbf{F}\) describes local stretching, compression, rotation, and
shearing.

%%%%%%%%%%%%%%%%%%%%%%%%%%%%%%%%%%%%%%%%%%%%%%%%%%%%%%%%%%%%
\subsection{Jacobian of the Deformation}
\label{subsec:deformation-jacobian}
%%%%%%%%%%%%%%%%%%%%%%%%%%%%%%%%%%%%%%%%%%%%%%%%%%%%%%%%%%%%

Define

\[
\boxed{
J
=
\det\mathbf{F}.
}
\]

The Jacobian measures local volume change.

If

\[
dV_0
\]

is a reference volume element, then

\[
\boxed{
dV
=
J\,dV_0.
}
\]

For physically admissible orientation-preserving deformations, one usually
requires

\[
\boxed{
J>0.
}
\]

%%%%%%%%%%%%%%%%%%%%%%%%%%%%%%%%%%%%%%%%%%%%%%%%%%%%%%%%%%%%
\subsection{Incompressibility}
\label{subsec:continuum-incompressibility}
%%%%%%%%%%%%%%%%%%%%%%%%%%%%%%%%%%%%%%%%%%%%%%%%%%%%%%%%%%%%

A material is incompressible when local volume is preserved.

In finite deformation,

\[
\boxed{
J=1.
}
\]

For an incompressible velocity field in the Eulerian description,

\[
\boxed{
\nabla\cdot\mathbf{v}=0.
}
\]

This condition will become central in fluid dynamics.

%%%%%%%%%%%%%%%%%%%%%%%%%%%%%%%%%%%%%%%%%%%%%%%%%%%%%%%%%%%%
\subsection{Right Cauchy--Green Tensor}
\label{subsec:right-cauchy-green}
%%%%%%%%%%%%%%%%%%%%%%%%%%%%%%%%%%%%%%%%%%%%%%%%%%%%%%%%%%%%

Define the right Cauchy--Green deformation tensor

\[
\boxed{
\mathbf{C}
=
\mathbf{F}^{T}\mathbf{F}.
}
\]

For a material line element,

\[
d\mathbf{x}\cdot d\mathbf{x}
=
d\mathbf{X}^{T}
\mathbf{C}
d\mathbf{X}.
\]

Thus \(\mathbf{C}\) measures changes in lengths and angles without including
rigid-body rotation.

%%%%%%%%%%%%%%%%%%%%%%%%%%%%%%%%%%%%%%%%%%%%%%%%%%%%%%%%%%%%
\subsection{Left Cauchy--Green Tensor}
\label{subsec:left-cauchy-green}
%%%%%%%%%%%%%%%%%%%%%%%%%%%%%%%%%%%%%%%%%%%%%%%%%%%%%%%%%%%%

The left Cauchy--Green tensor is

\[
\boxed{
\mathbf{B}
=
\mathbf{F}\mathbf{F}^{T}.
}
\]

Both \(\mathbf{B}\) and \(\mathbf{C}\) encode finite deformation, but they are
expressed relative to different configurations.

%%%%%%%%%%%%%%%%%%%%%%%%%%%%%%%%%%%%%%%%%%%%%%%%%%%%%%%%%%%%
\subsection{Green--Lagrange Strain}
\label{subsec:green-lagrange-strain}
%%%%%%%%%%%%%%%%%%%%%%%%%%%%%%%%%%%%%%%%%%%%%%%%%%%%%%%%%%%%

The Green--Lagrange strain tensor is

\[
\boxed{
\mathbf{E}
=
\frac12
\left(
\mathbf{C}-\mathbf{I}
\right).
}
\]

Equivalently,

\[
\boxed{
\mathbf{E}
=
\frac12
\left[
\nabla\mathbf{u}
+
(\nabla\mathbf{u})^T
+
(\nabla\mathbf{u})^T\nabla\mathbf{u}
\right].
}
\]

This strain measure is useful for finite deformations.

%%%%%%%%%%%%%%%%%%%%%%%%%%%%%%%%%%%%%%%%%%%%%%%%%%%%%%%%%%%%
\subsection{Infinitesimal Strain}
\label{subsec:infinitesimal-strain}
%%%%%%%%%%%%%%%%%%%%%%%%%%%%%%%%%%%%%%%%%%%%%%%%%%%%%%%%%%%%

When displacement gradients are small,

\[
\left\|
\nabla\mathbf{u}
\right\|
\ll1,
\]

quadratic terms may be neglected.

The infinitesimal strain tensor is

\[
\boxed{
\boldsymbol{\varepsilon}
=
\frac12
\left(
\nabla\mathbf{u}
+
(\nabla\mathbf{u})^T
\right).
}
\]

In components,

\[
\boxed{
\varepsilon_{ij}
=
\frac12
\left(
\frac{\partial u_i}{\partial x_j}
+
\frac{\partial u_j}{\partial x_i}
\right).
}
\]

%%%%%%%%%%%%%%%%%%%%%%%%%%%%%%%%%%%%%%%%%%%%%%%%%%%%%%%%%%%%
\subsection{Normal Strain}
\label{subsec:normal-strain}
%%%%%%%%%%%%%%%%%%%%%%%%%%%%%%%%%%%%%%%%%%%%%%%%%%%%%%%%%%%%

Diagonal strain components

\[
\boxed{
\varepsilon_{11},
\quad
\varepsilon_{22},
\quad
\varepsilon_{33}
}
\]

represent normal extension or compression along coordinate directions.

Positive normal strain corresponds to extension.

Negative normal strain corresponds to compression.

%%%%%%%%%%%%%%%%%%%%%%%%%%%%%%%%%%%%%%%%%%%%%%%%%%%%%%%%%%%%
\subsection{Shear Strain}
\label{subsec:shear-strain}
%%%%%%%%%%%%%%%%%%%%%%%%%%%%%%%%%%%%%%%%%%%%%%%%%%%%%%%%%%%%

Off-diagonal strain components

\[
\boxed{
\varepsilon_{12},
\quad
\varepsilon_{13},
\quad
\varepsilon_{23}
}
\]

measure changes in angles between material directions.

Engineering shear strain is often written

\[
\boxed{
\gamma_{12}
=
2\varepsilon_{12}.
}
\]

%%%%%%%%%%%%%%%%%%%%%%%%%%%%%%%%%%%%%%%%%%%%%%%%%%%%%%%%%%%%
\subsection{Rigid-Body Motion}
\label{subsec:rigid-body-motion-continuum}
%%%%%%%%%%%%%%%%%%%%%%%%%%%%%%%%%%%%%%%%%%%%%%%%%%%%%%%%%%%%

A rigid-body motion preserves distances.

A general rigid motion has form

\[
\boxed{
\mathbf{x}
=
\mathbf{Q}(t)\mathbf{X}
+
\mathbf{c}(t),
}
\]

where \(\mathbf{Q}\) is orthogonal:

\[
\boxed{
\mathbf{Q}^{T}\mathbf{Q}
=
\mathbf{I}.
}
\]

For such a motion,

\[
\mathbf{C}
=
\mathbf{I},
\]

so

\[
\boxed{
\mathbf{E}=0.
}
\]

Thus strain correctly vanishes under pure rigid-body motion.

%%%%%%%%%%%%%%%%%%%%%%%%%%%%%%%%%%%%%%%%%%%%%%%%%%%%%%%%%%%%
\subsection{Rate of Deformation}
\label{subsec:rate-of-deformation}
%%%%%%%%%%%%%%%%%%%%%%%%%%%%%%%%%%%%%%%%%%%%%%%%%%%%%%%%%%%%

The velocity gradient is

\[
\boxed{
\mathbf{L}
=
\nabla\mathbf{v}.
}
\]

It can be decomposed as

\[
\boxed{
\mathbf{L}
=
\mathbf{D}
+
\mathbf{W},
}
\]

where

\[
\boxed{
\mathbf{D}
=
\frac12
\left(
\nabla\mathbf{v}
+
(\nabla\mathbf{v})^T
\right)
}
\]

is the rate-of-deformation tensor and

\[
\boxed{
\mathbf{W}
=
\frac12
\left(
\nabla\mathbf{v}
-
(\nabla\mathbf{v})^T
\right)
}
\]

is the spin tensor.

%%%%%%%%%%%%%%%%%%%%%%%%%%%%%%%%%%%%%%%%%%%%%%%%%%%%%%%%%%%%
\subsection{Mass Density}
\label{subsec:continuum-density}
%%%%%%%%%%%%%%%%%%%%%%%%%%%%%%%%%%%%%%%%%%%%%%%%%%%%%%%%%%%%

Mass density is

\[
\boxed{
\rho
=
\frac{\text{mass}}
{\text{volume}}.
}
\]

The total mass of a spatial region \(\Omega_t\) is

\[
\boxed{
M
=
\int_{\Omega_t}
\rho(\mathbf{x},t)\,dV.
}
\]

For a material body, total mass is conserved.

%%%%%%%%%%%%%%%%%%%%%%%%%%%%%%%%%%%%%%%%%%%%%%%%%%%%%%%%%%%%
\subsection{Mass Conservation}
\label{subsec:continuum-mass-conservation}
%%%%%%%%%%%%%%%%%%%%%%%%%%%%%%%%%%%%%%%%%%%%%%%%%%%%%%%%%%%%

For a material region,

\[
\boxed{
\frac{d}{dt}
\int_{\Omega_t}
\rho\,dV
=
0.
}
\]

Using transport identities leads to the local continuity equation

\[
\boxed{
\frac{\partial\rho}{\partial t}
+
\nabla\cdot(\rho\mathbf{v})
=
0.
}
\]

%%%%%%%%%%%%%%%%%%%%%%%%%%%%%%%%%%%%%%%%%%%%%%%%%%%%%%%%%%%%
\subsection{Material Form of Mass Conservation}
\label{subsec:material-mass-conservation}
%%%%%%%%%%%%%%%%%%%%%%%%%%%%%%%%%%%%%%%%%%%%%%%%%%%%%%%%%%%%

Expanding,

\[
\nabla\cdot(\rho\mathbf{v})
=
\mathbf{v}\cdot\nabla\rho
+
\rho\nabla\cdot\mathbf{v}.
\]

Therefore,

\[
\boxed{
\frac{D\rho}{Dt}
+
\rho\nabla\cdot\mathbf{v}
=
0.
}
\]

For incompressible constant-density motion,

\[
\boxed{
\nabla\cdot\mathbf{v}=0.
}
\]

%%%%%%%%%%%%%%%%%%%%%%%%%%%%%%%%%%%%%%%%%%%%%%%%%%%%%%%%%%%%
\subsection{Reference Density Relation}
\label{subsec:reference-density-relation}
%%%%%%%%%%%%%%%%%%%%%%%%%%%%%%%%%%%%%%%%%%%%%%%%%%%%%%%%%%%%

Let \(\rho_0(\mathbf{X})\) be reference density.

Mass conservation gives

\[
\rho\,dV
=
\rho_0\,dV_0.
\]

Since

\[
dV
=
J\,dV_0,
\]

we obtain

\[
\boxed{
\rho J
=
\rho_0.
}
\]

Thus

\[
\boxed{
\rho
=
\frac{\rho_0}{J}.
}
\]

%%%%%%%%%%%%%%%%%%%%%%%%%%%%%%%%%%%%%%%%%%%%%%%%%%%%%%%%%%%%
\subsection{Body Forces}
\label{subsec:body-forces}
%%%%%%%%%%%%%%%%%%%%%%%%%%%%%%%%%%%%%%%%%%%%%%%%%%%%%%%%%%%%

Body forces act throughout the volume of a material.

Examples include:

\[
\boxed{
\text{gravity},
}
\]

\[
\boxed{
\text{electromagnetic forces},
}
\]

and

\[
\boxed{
\text{inertial forces in noninertial frames}.
}
\]

If \(\mathbf{b}\) is body force per unit mass, the force density is

\[
\boxed{
\rho\mathbf{b}.
}
\]

%%%%%%%%%%%%%%%%%%%%%%%%%%%%%%%%%%%%%%%%%%%%%%%%%%%%%%%%%%%%
\subsection{Surface Forces}
\label{subsec:surface-forces}
%%%%%%%%%%%%%%%%%%%%%%%%%%%%%%%%%%%%%%%%%%%%%%%%%%%%%%%%%%%%

Surface forces act across internal or external surfaces.

Define traction

\[
\boxed{
\mathbf{t}(\mathbf{n}),
}
\]

where \(\mathbf{n}\) is the unit normal to the surface.

Traction has units of force per unit area.

%%%%%%%%%%%%%%%%%%%%%%%%%%%%%%%%%%%%%%%%%%%%%%%%%%%%%%%%%%%%
\subsection{Cauchy's Stress Theorem}
\label{subsec:cauchy-stress-theorem}
%%%%%%%%%%%%%%%%%%%%%%%%%%%%%%%%%%%%%%%%%%%%%%%%%%%%%%%%%%%%

Cauchy's stress theorem states that traction depends linearly on the surface
normal:

\[
\boxed{
\mathbf{t}(\mathbf{n})
=
\boldsymbol{\sigma}\mathbf{n}.
}
\]

The tensor

\[
\boxed{
\boldsymbol{\sigma}
}
\]

is the Cauchy stress tensor.

%%%%%%%%%%%%%%%%%%%%%%%%%%%%%%%%%%%%%%%%%%%%%%%%%%%%%%%%%%%%
\subsection{Cauchy Stress Tensor}
\label{subsec:cauchy-stress-tensor}
%%%%%%%%%%%%%%%%%%%%%%%%%%%%%%%%%%%%%%%%%%%%%%%%%%%%%%%%%%%%

In Cartesian coordinates,

\[
\boxed{
\boldsymbol{\sigma}
=
\begin{bmatrix}
\sigma_{11} & \sigma_{12} & \sigma_{13}\\
\sigma_{21} & \sigma_{22} & \sigma_{23}\\
\sigma_{31} & \sigma_{32} & \sigma_{33}
\end{bmatrix}.
}
\]

Diagonal components represent normal stresses.

Off-diagonal components represent shear stresses.

%%%%%%%%%%%%%%%%%%%%%%%%%%%%%%%%%%%%%%%%%%%%%%%%%%%%%%%%%%%%
\subsection{Traction on a Plane}
\label{subsec:traction-plane}
%%%%%%%%%%%%%%%%%%%%%%%%%%%%%%%%%%%%%%%%%%%%%%%%%%%%%%%%%%%%

For unit normal

\[
\mathbf{n},
\]

the traction vector is

\[
\boxed{
\mathbf{t}
=
\boldsymbol{\sigma}\mathbf{n}.
}
\]

The normal stress is

\[
\boxed{
\sigma_n
=
\mathbf{n}\cdot
\boldsymbol{\sigma}\mathbf{n}.
}
\]

The remaining component is shear traction.

%%%%%%%%%%%%%%%%%%%%%%%%%%%%%%%%%%%%%%%%%%%%%%%%%%%%%%%%%%%%
\subsection{Worked Example: Traction Vector}
\label{subsec:worked-traction-vector}
%%%%%%%%%%%%%%%%%%%%%%%%%%%%%%%%%%%%%%%%%%%%%%%%%%%%%%%%%%%%

\begin{workedexamplebox}{Stress Acting on a Coordinate Plane}

Let

\[
\boldsymbol{\sigma}
=
\begin{bmatrix}
5 & 2 & 0\\
2 & 4 & 1\\
0 & 1 & 3
\end{bmatrix}
\]

and choose

\[
\mathbf{n}
=
\begin{bmatrix}
1\\
0\\
0
\end{bmatrix}.
\]

Then

\[
\mathbf{t}
=
\boldsymbol{\sigma}\mathbf{n}.
\]

Therefore,

\begin{answerbox}
\[
\boxed{
\mathbf{t}
=
\begin{bmatrix}
5\\
2\\
0
\end{bmatrix}.
}
\]
\end{answerbox}

The normal stress is \(5\), while the shear component has magnitude \(2\).

\end{workedexamplebox}

%%%%%%%%%%%%%%%%%%%%%%%%%%%%%%%%%%%%%%%%%%%%%%%%%%%%%%%%%%%%
\subsection{Balance of Linear Momentum}
\label{subsec:balance-linear-momentum}
%%%%%%%%%%%%%%%%%%%%%%%%%%%%%%%%%%%%%%%%%%%%%%%%%%%%%%%%%%%%

Newton's second law for a continuum gives

\[
\boxed{
\rho
\frac{D\mathbf{v}}{Dt}
=
\nabla\cdot\boldsymbol{\sigma}
+
\rho\mathbf{b}.
}
\]

In components,

\[
\boxed{
\rho
\frac{Dv_i}{Dt}
=
\frac{\partial\sigma_{ij}}
{\partial x_j}
+
\rho b_i.
}
\]

This is the local momentum balance.

%%%%%%%%%%%%%%%%%%%%%%%%%%%%%%%%%%%%%%%%%%%%%%%%%%%%%%%%%%%%
\subsection{Static Equilibrium}
\label{subsec:continuum-static-equilibrium}
%%%%%%%%%%%%%%%%%%%%%%%%%%%%%%%%%%%%%%%%%%%%%%%%%%%%%%%%%%%%

For a stationary body,

\[
\mathbf{v}=0,
\qquad
\mathbf{a}=0.
\]

Momentum balance becomes

\[
\boxed{
\nabla\cdot\boldsymbol{\sigma}
+
\rho\mathbf{b}
=
0.
}
\]

This is the equilibrium equation of solid mechanics.

%%%%%%%%%%%%%%%%%%%%%%%%%%%%%%%%%%%%%%%%%%%%%%%%%%%%%%%%%%%%
\subsection{Balance of Angular Momentum}
\label{subsec:balance-angular-momentum}
%%%%%%%%%%%%%%%%%%%%%%%%%%%%%%%%%%%%%%%%%%%%%%%%%%%%%%%%%%%%

In a classical continuum without couple stresses, balance of angular
momentum implies

\[
\boxed{
\boldsymbol{\sigma}
=
\boldsymbol{\sigma}^{T}.
}
\]

Thus

\[
\boxed{
\sigma_{ij}
=
\sigma_{ji}.
}
\]

The Cauchy stress tensor is symmetric.

%%%%%%%%%%%%%%%%%%%%%%%%%%%%%%%%%%%%%%%%%%%%%%%%%%%%%%%%%%%%
\subsection{Principal Stresses}
\label{subsec:principal-stresses}
%%%%%%%%%%%%%%%%%%%%%%%%%%%%%%%%%%%%%%%%%%%%%%%%%%%%%%%%%%%%

A principal stress direction is a direction in which traction is purely
normal.

Thus

\[
\boxed{
\boldsymbol{\sigma}\mathbf{n}
=
\lambda\mathbf{n}.
}
\]

Principal stresses are therefore eigenvalues of the stress tensor.

Since \(\boldsymbol{\sigma}\) is symmetric, its principal stresses are real
and its principal directions can be chosen orthogonal.

%%%%%%%%%%%%%%%%%%%%%%%%%%%%%%%%%%%%%%%%%%%%%%%%%%%%%%%%%%%%
\subsection{Worked Example: Principal Stress}
\label{subsec:worked-principal-stress}
%%%%%%%%%%%%%%%%%%%%%%%%%%%%%%%%%%%%%%%%%%%%%%%%%%%%%%%%%%%%

\begin{workedexamplebox}{Two-Dimensional Stress State}

Let

\[
\boldsymbol{\sigma}
=
\begin{bmatrix}
6 & 2\\
2 & 3
\end{bmatrix}.
\]

Principal stresses satisfy

\[
\det
\left(
\boldsymbol{\sigma}
-
\lambda\mathbf{I}
\right)
=
0.
\]

Thus

\[
\begin{vmatrix}
6-\lambda & 2\\
2 & 3-\lambda
\end{vmatrix}
=
0.
\]

Therefore,

\[
(6-\lambda)(3-\lambda)-4=0,
\]

so

\[
\lambda^2-9\lambda+14=0.
\]

Hence

\[
\begin{answerbox}
\[
\boxed{
\lambda_1=7,
\qquad
\lambda_2=2.
}
\]
\end{answerbox}

\end{workedexamplebox}

%%%%%%%%%%%%%%%%%%%%%%%%%%%%%%%%%%%%%%%%%%%%%%%%%%%%%%%%%%%%
\subsection{Hydrostatic Stress}
\label{subsec:hydrostatic-stress}
%%%%%%%%%%%%%%%%%%%%%%%%%%%%%%%%%%%%%%%%%%%%%%%%%%%%%%%%%%%%

A purely hydrostatic stress state has form

\[
\boxed{
\boldsymbol{\sigma}
=
-p\mathbf{I}.
}
\]

The scalar \(p\) is pressure.

Hydrostatic stress produces equal normal compression in every direction and
no shear stress.

%%%%%%%%%%%%%%%%%%%%%%%%%%%%%%%%%%%%%%%%%%%%%%%%%%%%%%%%%%%%
\subsection{Deviatoric Stress}
\label{subsec:deviatoric-stress}
%%%%%%%%%%%%%%%%%%%%%%%%%%%%%%%%%%%%%%%%%%%%%%%%%%%%%%%%%%%%

The stress tensor can be decomposed into hydrostatic and deviatoric parts:

\[
\boxed{
\boldsymbol{\sigma}
=
-p\mathbf{I}
+
\mathbf{s}.
}
\]

The deviatoric stress satisfies

\[
\boxed{
\operatorname{tr}\mathbf{s}
=
0.
}
\]

The deviatoric part is associated with shape-changing deformation.

%%%%%%%%%%%%%%%%%%%%%%%%%%%%%%%%%%%%%%%%%%%%%%%%%%%%%%%%%%%%
\subsection{Constitutive Equations}
\label{subsec:constitutive-equations}
%%%%%%%%%%%%%%%%%%%%%%%%%%%%%%%%%%%%%%%%%%%%%%%%%%%%%%%%%%%%

Balance laws alone do not determine material behavior.

A constitutive equation specifies how stress depends on deformation,
deformation rate, temperature, or other variables.

Schematically,

\[
\boxed{
\boldsymbol{\sigma}
=
\mathcal{F}
(
\boldsymbol{\varepsilon},
\dot{\boldsymbol{\varepsilon}},
T,\ldots
).
}
\]

Different materials require different constitutive laws.

%%%%%%%%%%%%%%%%%%%%%%%%%%%%%%%%%%%%%%%%%%%%%%%%%%%%%%%%%%%%
\subsection{Three Classes of Equations}
\label{subsec:three-continuum-equation-classes}
%%%%%%%%%%%%%%%%%%%%%%%%%%%%%%%%%%%%%%%%%%%%%%%%%%%%%%%%%%%%

Continuum models typically contain:

\[
\boxed{
\text{kinematic equations},
}
\]

which describe deformation;

\[
\boxed{
\text{balance laws},
}
\]

which express conservation;

and

\[
\boxed{
\text{constitutive equations},
}
\]

which describe material response.

All three are needed for a closed physical model.

%%%%%%%%%%%%%%%%%%%%%%%%%%%%%%%%%%%%%%%%%%%%%%%%%%%%%%%%%%%%
\subsection{Linear Elasticity}
\label{subsec:linear-elasticity}
%%%%%%%%%%%%%%%%%%%%%%%%%%%%%%%%%%%%%%%%%%%%%%%%%%%%%%%%%%%%

Linear elasticity assumes:

\[
\boxed{
\text{small deformation}
}
\]

and

\[
\boxed{
\text{linear stress-strain response}.
}
\]

The constitutive law is

\[
\boxed{
\sigma_{ij}
=
C_{ijkl}
\varepsilon_{kl},
}
\]

where \(C_{ijkl}\) is the elasticity tensor.

%%%%%%%%%%%%%%%%%%%%%%%%%%%%%%%%%%%%%%%%%%%%%%%%%%%%%%%%%%%%
\subsection{Isotropic Linear Elasticity}
\label{subsec:isotropic-linear-elasticity}
%%%%%%%%%%%%%%%%%%%%%%%%%%%%%%%%%%%%%%%%%%%%%%%%%%%%%%%%%%%%

For an isotropic material,

\[
\boxed{
\boldsymbol{\sigma}
=
\lambda
\operatorname{tr}
(\boldsymbol{\varepsilon})
\mathbf{I}
+
2\mu
\boldsymbol{\varepsilon}.
}
\]

The constants

\[
\lambda,
\qquad
\mu
\]

are the Lamé parameters.

The parameter \(\mu\) is also the shear modulus.

%%%%%%%%%%%%%%%%%%%%%%%%%%%%%%%%%%%%%%%%%%%%%%%%%%%%%%%%%%%%
\subsection{Young's Modulus and Poisson Ratio}
\label{subsec:young-poisson}
%%%%%%%%%%%%%%%%%%%%%%%%%%%%%%%%%%%%%%%%%%%%%%%%%%%%%%%%%%%%

Isotropic elastic behavior may instead be described by Young's modulus \(E\)
and Poisson ratio \(\nu\).

The relationships include

\[
\boxed{
\mu
=
\frac{E}
{2(1+\nu)},
}
\]

and

\[
\boxed{
\lambda
=
\frac{E\nu}
{(1+\nu)(1-2\nu)}.
}
\]

For stable ordinary isotropic elastic materials, parameters satisfy suitable
positivity restrictions.

%%%%%%%%%%%%%%%%%%%%%%%%%%%%%%%%%%%%%%%%%%%%%%%%%%%%%%%%%%%%
\subsection{Uniaxial Stress}
\label{subsec:uniaxial-stress}
%%%%%%%%%%%%%%%%%%%%%%%%%%%%%%%%%%%%%%%%%%%%%%%%%%%%%%%%%%%%

For simple one-dimensional linear elasticity,

\[
\boxed{
\sigma
=
E\varepsilon.
}
\]

If a bar of length \(L\) changes by \(\Delta L\), then

\[
\boxed{
\varepsilon
=
\frac{\Delta L}{L}.
}
\]

Thus

\[
\boxed{
\sigma
=
E\frac{\Delta L}{L}.
}
\]

%%%%%%%%%%%%%%%%%%%%%%%%%%%%%%%%%%%%%%%%%%%%%%%%%%%%%%%%%%%%
\subsection{Worked Example: Elastic Bar}
\label{subsec:worked-elastic-bar}
%%%%%%%%%%%%%%%%%%%%%%%%%%%%%%%%%%%%%%%%%%%%%%%%%%%%%%%%%%%%

\begin{workedexamplebox}{Extension of a Uniform Bar}

Suppose a bar has

\[
E=200\text{ GPa},
\]

length

\[
L=2\text{ m},
\]

and experiences uniform stress

\[
\sigma=100\text{ MPa}.
\]

From Hooke's law,

\[
\varepsilon
=
\frac{\sigma}{E}.
\]

Thus

\[
\varepsilon
=
\frac{100\times10^6}
{200\times10^9}
=
5\times10^{-4}.
\]

The extension is

\[
\Delta L
=
\varepsilon L.
\]

Hence

\[
\begin{answerbox}
\[
\boxed{
\Delta L
=
10^{-3}\text{ m}
=
1\text{ mm}.
}
\]
\end{answerbox}

\end{workedexamplebox}

%%%%%%%%%%%%%%%%%%%%%%%%%%%%%%%%%%%%%%%%%%%%%%%%%%%%%%%%%%%%
\subsection{Navier--Cauchy Equations}
\label{subsec:navier-cauchy-equations}
%%%%%%%%%%%%%%%%%%%%%%%%%%%%%%%%%%%%%%%%%%%%%%%%%%%%%%%%%%%%

For homogeneous isotropic linear elasticity,

\[
\boldsymbol{\sigma}
=
\lambda
(\nabla\cdot\mathbf{u})
\mathbf{I}
+
\mu
\left(
\nabla\mathbf{u}
+
(\nabla\mathbf{u})^T
\right).
\]

Substituting into momentum balance gives

\[
\boxed{
\rho\mathbf{u}_{tt}
=
(\lambda+\mu)
\nabla(\nabla\cdot\mathbf{u})
+
\mu\nabla^2\mathbf{u}
+
\rho\mathbf{b}.
}
\]

These are the Navier--Cauchy equations of linear elastodynamics.

%%%%%%%%%%%%%%%%%%%%%%%%%%%%%%%%%%%%%%%%%%%%%%%%%%%%%%%%%%%%
\subsection{Static Linear Elasticity}
\label{subsec:static-linear-elasticity}
%%%%%%%%%%%%%%%%%%%%%%%%%%%%%%%%%%%%%%%%%%%%%%%%%%%%%%%%%%%%

For static equilibrium,

\[
\mathbf{u}_{tt}=0.
\]

Then

\[
\boxed{
(\lambda+\mu)
\nabla(\nabla\cdot\mathbf{u})
+
\mu\nabla^2\mathbf{u}
+
\rho\mathbf{b}
=
0.
}
\]

Boundary conditions are required to determine a unique displacement field.

%%%%%%%%%%%%%%%%%%%%%%%%%%%%%%%%%%%%%%%%%%%%%%%%%%%%%%%%%%%%
\subsection{Elastic Strain Energy}
\label{subsec:elastic-strain-energy}
%%%%%%%%%%%%%%%%%%%%%%%%%%%%%%%%%%%%%%%%%%%%%%%%%%%%%%%%%%%%

For linear elasticity, strain-energy density is

\[
\boxed{
W
=
\frac12
\boldsymbol{\sigma}
:
\boldsymbol{\varepsilon}.
}
\]

The colon denotes double contraction:

\[
\boxed{
\mathbf{A}:\mathbf{B}
=
A_{ij}B_{ij}.
}
\]

For one-dimensional elasticity,

\[
\boxed{
W
=
\frac12E\varepsilon^2.
}
\]

%%%%%%%%%%%%%%%%%%%%%%%%%%%%%%%%%%%%%%%%%%%%%%%%%%%%%%%%%%%%
\subsection{Hyperelasticity}
\label{subsec:hyperelasticity}
%%%%%%%%%%%%%%%%%%%%%%%%%%%%%%%%%%%%%%%%%%%%%%%%%%%%%%%%%%%%

A hyperelastic material is described by a strain-energy function

\[
\boxed{
W=W(\mathbf{F})
}
\]

or equivalently by objective strain measures.

Stress is obtained by differentiating the energy with respect to the
appropriate deformation measure.

Hyperelastic models are useful for large elastic deformations.

%%%%%%%%%%%%%%%%%%%%%%%%%%%%%%%%%%%%%%%%%%%%%%%%%%%%%%%%%%%%
\subsection{Objectivity}
\label{subsec:continuum-objectivity}
%%%%%%%%%%%%%%%%%%%%%%%%%%%%%%%%%%%%%%%%%%%%%%%%%%%%%%%%%%%%

A physically acceptable constitutive law should not depend on arbitrary
rigid rotations of the observer.

This requirement is called frame indifference or objectivity.

It ensures that material response depends on physical deformation rather
than coordinate orientation.

%%%%%%%%%%%%%%%%%%%%%%%%%%%%%%%%%%%%%%%%%%%%%%%%%%%%%%%%%%%%
\subsection{Viscoelasticity}
\label{subsec:viscoelasticity}
%%%%%%%%%%%%%%%%%%%%%%%%%%%%%%%%%%%%%%%%%%%%%%%%%%%%%%%%%%%%

Viscoelastic materials exhibit both elastic and viscous behavior.

Stress can depend on deformation and deformation rate.

A one-dimensional constitutive relation may include both

\[
\boxed{
\varepsilon
}
\]

and

\[
\boxed{
\dot{\varepsilon}.
}
\]

Such materials display time-dependent relaxation and creep.

%%%%%%%%%%%%%%%%%%%%%%%%%%%%%%%%%%%%%%%%%%%%%%%%%%%%%%%%%%%%
\subsection{Kelvin--Voigt Model}
\label{subsec:kelvin-voigt-model}
%%%%%%%%%%%%%%%%%%%%%%%%%%%%%%%%%%%%%%%%%%%%%%%%%%%%%%%%%%%%

The Kelvin--Voigt constitutive equation is

\[
\boxed{
\sigma
=
E\varepsilon
+
\eta\dot{\varepsilon}.
}
\]

The spring and dashpot act in parallel.

The model captures delayed deformation under applied load.

%%%%%%%%%%%%%%%%%%%%%%%%%%%%%%%%%%%%%%%%%%%%%%%%%%%%%%%%%%%%
\subsection{Maxwell Model}
\label{subsec:maxwell-viscoelastic-model}
%%%%%%%%%%%%%%%%%%%%%%%%%%%%%%%%%%%%%%%%%%%%%%%%%%%%%%%%%%%%

For a Maxwell material,

\[
\boxed{
\dot{\varepsilon}
=
\frac{1}{E}\dot{\sigma}
+
\frac{\sigma}{\eta}.
}
\]

A spring and dashpot act in series.

The model can exhibit stress relaxation under fixed deformation.

%%%%%%%%%%%%%%%%%%%%%%%%%%%%%%%%%%%%%%%%%%%%%%%%%%%%%%%%%%%%
\subsection{Worked Example: Kelvin--Voigt Creep}
\label{subsec:worked-kelvin-voigt-creep}
%%%%%%%%%%%%%%%%%%%%%%%%%%%%%%%%%%%%%%%%%%%%%%%%%%%%%%%%%%%%

\begin{workedexamplebox}{Constant Applied Stress}

Suppose

\[
\sigma(t)=\sigma_0
\]

and the Kelvin--Voigt law is

\[
\sigma_0
=
E\varepsilon
+
\eta\dot{\varepsilon}.
\]

Then

\[
\eta\dot{\varepsilon}
+
E\varepsilon
=
\sigma_0.
\]

With

\[
\varepsilon(0)=0,
\]

the solution is

\[
\begin{answerbox}
\[
\boxed{
\varepsilon(t)
=
\frac{\sigma_0}{E}
\left(
1-e^{-Et/\eta}
\right).
}
\]
\end{answerbox}

The strain approaches the elastic value \(\sigma_0/E\) gradually.

\end{workedexamplebox}

%%%%%%%%%%%%%%%%%%%%%%%%%%%%%%%%%%%%%%%%%%%%%%%%%%%%%%%%%%%%
\subsection{Plasticity}
\label{subsec:continuum-plasticity}
%%%%%%%%%%%%%%%%%%%%%%%%%%%%%%%%%%%%%%%%%%%%%%%%%%%%%%%%%%%%

Plasticity describes irreversible deformation.

A material may behave elastically while stress remains below a yield
threshold.

After yielding, permanent deformation can remain even when loading is
removed.

A typical decomposition is

\[
\boxed{
\boldsymbol{\varepsilon}
=
\boldsymbol{\varepsilon}^{e}
+
\boldsymbol{\varepsilon}^{p},
}
\]

where superscripts \(e\) and \(p\) denote elastic and plastic parts.

%%%%%%%%%%%%%%%%%%%%%%%%%%%%%%%%%%%%%%%%%%%%%%%%%%%%%%%%%%%%
\subsection{Yield Criteria}
\label{subsec:yield-criteria}
%%%%%%%%%%%%%%%%%%%%%%%%%%%%%%%%%%%%%%%%%%%%%%%%%%%%%%%%%%%%

A yield criterion determines when plastic deformation begins.

Conceptually,

\[
\boxed{
f(\boldsymbol{\sigma})
\leq0
}
\]

describes elastic states.

Yield occurs when

\[
\boxed{
f(\boldsymbol{\sigma})=0.
}
\]

Common criteria are based on principal or deviatoric stresses.

%%%%%%%%%%%%%%%%%%%%%%%%%%%%%%%%%%%%%%%%%%%%%%%%%%%%%%%%%%%%
\subsection{Thermal Strain}
\label{subsec:thermal-strain}
%%%%%%%%%%%%%%%%%%%%%%%%%%%%%%%%%%%%%%%%%%%%%%%%%%%%%%%%%%%%

Temperature changes can produce deformation.

In one dimension,

\[
\boxed{
\varepsilon_{\text{thermal}}
=
\alpha
\Delta T,
}
\]

where

\[
\alpha
\]

is the coefficient of thermal expansion.

The total strain may be decomposed into mechanical and thermal contributions.

%%%%%%%%%%%%%%%%%%%%%%%%%%%%%%%%%%%%%%%%%%%%%%%%%%%%%%%%%%%%
\subsection{Thermoelasticity}
\label{subsec:thermoelasticity}
%%%%%%%%%%%%%%%%%%%%%%%%%%%%%%%%%%%%%%%%%%%%%%%%%%%%%%%%%%%%

For isotropic linear thermoelasticity, one may write schematically

\[
\boxed{
\boldsymbol{\sigma}
=
\mathsf{C}
:
\left(
\boldsymbol{\varepsilon}
-
\boldsymbol{\varepsilon}_{\text{thermal}}
\right).
}
\]

Thus stress depends on deformation relative to the material's preferred
thermally expanded configuration.

%%%%%%%%%%%%%%%%%%%%%%%%%%%%%%%%%%%%%%%%%%%%%%%%%%%%%%%%%%%%
\subsection{Energy Balance}
\label{subsec:continuum-energy-balance}
%%%%%%%%%%%%%%%%%%%%%%%%%%%%%%%%%%%%%%%%%%%%%%%%%%%%%%%%%%%%

A continuum also satisfies conservation of energy.

A local form can be written schematically as

\[
\boxed{
\rho\frac{De}{Dt}
=
\boldsymbol{\sigma}:\mathbf{D}
-
\nabla\cdot\mathbf{q}
+
\rho r,
}
\]

where

\[
e
=
\text{internal energy per unit mass},
\]

\[
\mathbf{q}
=
\text{heat flux},
\]

and

\[
r
=
\text{heat supply per unit mass}.
\]

%%%%%%%%%%%%%%%%%%%%%%%%%%%%%%%%%%%%%%%%%%%%%%%%%%%%%%%%%%%%
\subsection{Fourier Heat Conduction}
\label{subsec:fourier-heat-conduction}
%%%%%%%%%%%%%%%%%%%%%%%%%%%%%%%%%%%%%%%%%%%%%%%%%%%%%%%%%%%%

Fourier's constitutive law is

\[
\boxed{
\mathbf{q}
=
-k\nabla T,
}
\]

where \(k>0\) is thermal conductivity.

Heat flows down the temperature gradient.

Combining this with energy balance leads to heat-conduction equations.

%%%%%%%%%%%%%%%%%%%%%%%%%%%%%%%%%%%%%%%%%%%%%%%%%%%%%%%%%%%%
\subsection{Second Law of Thermodynamics}
\label{subsec:continuum-second-law}
%%%%%%%%%%%%%%%%%%%%%%%%%%%%%%%%%%%%%%%%%%%%%%%%%%%%%%%%%%%%

The second law places restrictions on constitutive models.

It requires nonnegative entropy production.

Conceptually,

\[
\boxed{
\text{constitutive equations cannot be chosen arbitrarily}.
}
\]

They must satisfy both mechanical balance laws and thermodynamic
admissibility.

%%%%%%%%%%%%%%%%%%%%%%%%%%%%%%%%%%%%%%%%%%%%%%%%%%%%%%%%%%%%
\subsection{Boundary Conditions in Solid Mechanics}
\label{subsec:solid-boundary-conditions}
%%%%%%%%%%%%%%%%%%%%%%%%%%%%%%%%%%%%%%%%%%%%%%%%%%%%%%%%%%%%

Two common boundary conditions are displacement and traction conditions.

On one part of the boundary,

\[
\boxed{
\mathbf{u}
=
\bar{\mathbf{u}}.
}
\]

On another part,

\[
\boxed{
\boldsymbol{\sigma}\mathbf{n}
=
\bar{\mathbf{t}}.
}
\]

These are analogous to Dirichlet and Neumann boundary conditions.

%%%%%%%%%%%%%%%%%%%%%%%%%%%%%%%%%%%%%%%%%%%%%%%%%%%%%%%%%%%%
\subsection{Initial Conditions}
\label{subsec:continuum-initial-conditions}
%%%%%%%%%%%%%%%%%%%%%%%%%%%%%%%%%%%%%%%%%%%%%%%%%%%%%%%%%%%%

Dynamic solid mechanics typically requires

\[
\boxed{
\mathbf{u}(\mathbf{x},0)
=
\mathbf{u}_0(\mathbf{x}),
}
\]

and

\[
\boxed{
\mathbf{u}_t(\mathbf{x},0)
=
\mathbf{v}_0(\mathbf{x}).
}
\]

These specify the initial displacement and velocity fields.

%%%%%%%%%%%%%%%%%%%%%%%%%%%%%%%%%%%%%%%%%%%%%%%%%%%%%%%%%%%%
\subsection{Elastic Waves}
\label{subsec:elastic-waves}
%%%%%%%%%%%%%%%%%%%%%%%%%%%%%%%%%%%%%%%%%%%%%%%%%%%%%%%%%%%%

The Navier--Cauchy equations support wave propagation.

In an isotropic elastic material, two important wave types arise:

\[
\boxed{
\text{longitudinal waves}
}
\]

and

\[
\boxed{
\text{transverse waves}.
}
\]

Their speeds depend on elastic moduli and density.

%%%%%%%%%%%%%%%%%%%%%%%%%%%%%%%%%%%%%%%%%%%%%%%%%%%%%%%%%%%%
\subsection{Longitudinal Wave Speed}
\label{subsec:longitudinal-wave-speed}
%%%%%%%%%%%%%%%%%%%%%%%%%%%%%%%%%%%%%%%%%%%%%%%%%%%%%%%%%%%%

The longitudinal or pressure-wave speed is

\[
\boxed{
c_P
=
\sqrt{
\frac{\lambda+2\mu}
{\rho}
}.
}
\]

Particle displacement is parallel to the propagation direction.

These are also called \(P\)-waves.

%%%%%%%%%%%%%%%%%%%%%%%%%%%%%%%%%%%%%%%%%%%%%%%%%%%%%%%%%%%%
\subsection{Shear Wave Speed}
\label{subsec:shear-wave-speed}
%%%%%%%%%%%%%%%%%%%%%%%%%%%%%%%%%%%%%%%%%%%%%%%%%%%%%%%%%%%%

The transverse or shear-wave speed is

\[
\boxed{
c_S
=
\sqrt{
\frac{\mu}{\rho}
}.
}
\]

Particle displacement is perpendicular to the propagation direction.

These are also called \(S\)-waves.

%%%%%%%%%%%%%%%%%%%%%%%%%%%%%%%%%%%%%%%%%%%%%%%%%%%%%%%%%%%%
\subsection{Worked Example: Elastic Wave Speeds}
\label{subsec:worked-elastic-wave-speed}
%%%%%%%%%%%%%%%%%%%%%%%%%%%%%%%%%%%%%%%%%%%%%%%%%%%%%%%%%%%%

\begin{workedexamplebox}{Comparing \(P\)- and \(S\)-Wave Speeds}

Suppose

\[
\lambda=4,
\qquad
\mu=3,
\qquad
\rho=2.
\]

Then

\[
c_P
=
\sqrt{
\frac{4+2(3)}
{2}
}
=
\sqrt{5},
\]

while

\[
c_S
=
\sqrt{
\frac{3}{2}
}.
\]

Therefore,

\begin{answerbox}
\[
\boxed{
c_P=\sqrt{5},
\qquad
c_S=\sqrt{\frac32}.
}
\]
\end{answerbox}

Since \(c_P>c_S\), the longitudinal wave travels faster.

\end{workedexamplebox}

%%%%%%%%%%%%%%%%%%%%%%%%%%%%%%%%%%%%%%%%%%%%%%%%%%%%%%%%%%%%
\subsection{Weak Form of Elasticity}
\label{subsec:weak-form-elasticity}
%%%%%%%%%%%%%%%%%%%%%%%%%%%%%%%%%%%%%%%%%%%%%%%%%%%%%%%%%%%%

Consider static equilibrium

\[
-\nabla\cdot\boldsymbol{\sigma}
=
\mathbf{f}.
\]

Multiply by a test function \(\mathbf{w}\) and integrate:

\[
-\int_\Omega
\mathbf{w}\cdot
\nabla\cdot\boldsymbol{\sigma}
\,dV
=
\int_\Omega
\mathbf{w}\cdot\mathbf{f}
\,dV.
\]

Integration by parts gives a weak formulation involving

\[
\boxed{
\nabla\mathbf{w}
}
\]

rather than second derivatives of displacement.

This is the foundation of finite element formulations.

%%%%%%%%%%%%%%%%%%%%%%%%%%%%%%%%%%%%%%%%%%%%%%%%%%%%%%%%%%%%
\subsection{Principle of Virtual Work}
\label{subsec:virtual-work}
%%%%%%%%%%%%%%%%%%%%%%%%%%%%%%%%%%%%%%%%%%%%%%%%%%%%%%%%%%%%

A common weak equilibrium statement is

\[
\boxed{
\delta W_{\text{internal}}
=
\delta W_{\text{external}}.
}
\]

For linear elasticity,

\[
\boxed{
\int_\Omega
\boldsymbol{\varepsilon}(\delta\mathbf{u})
:
\boldsymbol{\sigma}
\,dV
=
\int_\Omega
\delta\mathbf{u}\cdot\mathbf{f}
\,dV
+
\int_{\partial\Omega_t}
\delta\mathbf{u}\cdot\bar{\mathbf{t}}
\,dA.
}
\]

This provides a variational formulation of equilibrium.

%%%%%%%%%%%%%%%%%%%%%%%%%%%%%%%%%%%%%%%%%%%%%%%%%%%%%%%%%%%%
\subsection{Finite Element Method Preview}
\label{subsec:finite-element-continuum-preview}
%%%%%%%%%%%%%%%%%%%%%%%%%%%%%%%%%%%%%%%%%%%%%%%%%%%%%%%%%%%%

Finite element methods approximate a continuum by dividing the domain into
smaller elements.

The displacement field is approximated as

\[
\boxed{
\mathbf{u}_h
=
\sum_{i=1}^{N}
N_i(\mathbf{x})\mathbf{u}_i,
}
\]

where \(N_i\) are shape functions.

The weak form leads to a matrix system

\[
\boxed{
K\mathbf{u}
=
\mathbf{f}.
}
\]

%%%%%%%%%%%%%%%%%%%%%%%%%%%%%%%%%%%%%%%%%%%%%%%%%%%%%%%%%%%%
\subsection{Stiffness Matrix}
\label{subsec:finite-element-stiffness}
%%%%%%%%%%%%%%%%%%%%%%%%%%%%%%%%%%%%%%%%%%%%%%%%%%%%%%%%%%%%

In linear elasticity, the finite element stiffness matrix has schematic form

\[
\boxed{
K
=
\int_\Omega
B^TDB
\,dV,
}
\]

where

\[
B
\]

maps nodal displacements to strains and

\[
D
\]

represents the constitutive elasticity matrix.

This connects continuum mechanics with numerical linear algebra.

%%%%%%%%%%%%%%%%%%%%%%%%%%%%%%%%%%%%%%%%%%%%%%%%%%%%%%%%%%%%
\subsection{Dimensional Analysis}
\label{subsec:continuum-dimensional-analysis}
%%%%%%%%%%%%%%%%%%%%%%%%%%%%%%%%%%%%%%%%%%%%%%%%%%%%%%%%%%%%

Important dimensions include:

\[
[\rho]
=
ML^{-3},
\]

\[
[\sigma]
=
ML^{-1}T^{-2},
\]

and

\[
[E]
=
ML^{-1}T^{-2}.
\]

Stress and elastic modulus therefore have the same physical dimensions.

Strain is dimensionless.

%%%%%%%%%%%%%%%%%%%%%%%%%%%%%%%%%%%%%%%%%%%%%%%%%%%%%%%%%%%%
\subsection{Characteristic Scales}
\label{subsec:continuum-characteristic-scales}
%%%%%%%%%%%%%%%%%%%%%%%%%%%%%%%%%%%%%%%%%%%%%%%%%%%%%%%%%%%%

Suppose characteristic quantities are:

\[
L
=
\text{length scale},
\]

\[
U
=
\text{velocity scale},
\]

\[
T
=
\text{time scale},
\]

and

\[
\Sigma
=
\text{stress scale}.
\]

Define dimensionless variables such as

\[
\boxed{
\hat{\mathbf{x}}
=
\frac{\mathbf{x}}{L},
}
\]

\[
\boxed{
\hat{\mathbf{v}}
=
\frac{\mathbf{v}}{U},
}
\]

and

\[
\boxed{
\hat{\boldsymbol{\sigma}}
=
\frac{\boldsymbol{\sigma}}{\Sigma}.
}
\]

Nondimensionalization reveals competing physical effects.

%%%%%%%%%%%%%%%%%%%%%%%%%%%%%%%%%%%%%%%%%%%%%%%%%%%%%%%%%%%%
\subsection{Small-Strain Approximation}
\label{subsec:small-strain-approximation}
%%%%%%%%%%%%%%%%%%%%%%%%%%%%%%%%%%%%%%%%%%%%%%%%%%%%%%%%%%%%

Linearized strain assumes

\[
\boxed{
\left\|
\nabla\mathbf{u}
\right\|
\ll1.
}
\]

Then quadratic terms in deformation may be neglected.

This approximation can fail when:

\[
\boxed{
\text{rotations are large},
}
\]

\[
\boxed{
\text{strains are large},
}
\]

or

\[
\boxed{
\text{geometry changes substantially}.
}
\]

%%%%%%%%%%%%%%%%%%%%%%%%%%%%%%%%%%%%%%%%%%%%%%%%%%%%%%%%%%%%
\subsection{Small Strain Versus Small Displacement}
\label{subsec:small-strain-small-displacement}
%%%%%%%%%%%%%%%%%%%%%%%%%%%%%%%%%%%%%%%%%%%%%%%%%%%%%%%%%%%%

Small strain does not always mean that absolute displacement is small.

A body can undergo a large rigid translation while experiencing zero strain.

Likewise, large rigid rotation can create difficulties for a linearized
description even though material distances remain unchanged.

The important quantity is the deformation gradient relative to rigid motion.

%%%%%%%%%%%%%%%%%%%%%%%%%%%%%%%%%%%%%%%%%%%%%%%%%%%%%%%%%%%%
\subsection{Finite Deformation Mechanics}
\label{subsec:finite-deformation-mechanics}
%%%%%%%%%%%%%%%%%%%%%%%%%%%%%%%%%%%%%%%%%%%%%%%%%%%%%%%%%%%%

When deformation is large, one must retain nonlinear terms.

Important quantities include

\[
\boxed{
\mathbf{F},
}
\]

\[
\boxed{
J,
}
\]

\[
\boxed{
\mathbf{C},
}
\]

and

\[
\boxed{
\mathbf{E}.
}
\]

Stress measures also require careful distinction between reference and
current configurations.

%%%%%%%%%%%%%%%%%%%%%%%%%%%%%%%%%%%%%%%%%%%%%%%%%%%%%%%%%%%%
\subsection{Stress Measures at Finite Deformation}
\label{subsec:finite-stress-measures}
%%%%%%%%%%%%%%%%%%%%%%%%%%%%%%%%%%%%%%%%%%%%%%%%%%%%%%%%%%%%

At finite deformation, several stress tensors are useful.

The Cauchy stress

\[
\boxed{
\boldsymbol{\sigma}
}
\]

acts in the current configuration.

The first Piola--Kirchhoff stress is often written

\[
\boxed{
\mathbf{P}.
}
\]

The second Piola--Kirchhoff stress is often written

\[
\boxed{
\mathbf{S}.
}
\]

Different stress measures are conjugate to different deformation measures.

%%%%%%%%%%%%%%%%%%%%%%%%%%%%%%%%%%%%%%%%%%%%%%%%%%%%%%%%%%%%
\subsection{First Piola--Kirchhoff Stress}
\label{subsec:first-piola-kirchhoff}
%%%%%%%%%%%%%%%%%%%%%%%%%%%%%%%%%%%%%%%%%%%%%%%%%%%%%%%%%%%%

The first Piola--Kirchhoff stress relates forces in the current
configuration to areas in the reference configuration.

It is related to Cauchy stress by

\[
\boxed{
\mathbf{P}
=
J
\boldsymbol{\sigma}
\mathbf{F}^{-T}.
}
\]

Unlike the Cauchy stress, \(\mathbf{P}\) need not be symmetric.

%%%%%%%%%%%%%%%%%%%%%%%%%%%%%%%%%%%%%%%%%%%%%%%%%%%%%%%%%%%%
\subsection{Second Piola--Kirchhoff Stress}
\label{subsec:second-piola-kirchhoff}
%%%%%%%%%%%%%%%%%%%%%%%%%%%%%%%%%%%%%%%%%%%%%%%%%%%%%%%%%%%%

The second Piola--Kirchhoff stress is

\[
\boxed{
\mathbf{S}
=
\mathbf{F}^{-1}\mathbf{P}.
}
\]

Equivalently,

\[
\boxed{
\mathbf{S}
=
J
\mathbf{F}^{-1}
\boldsymbol{\sigma}
\mathbf{F}^{-T}.
}
\]

It is useful with Green--Lagrange strain in reference-configuration
formulations.

%%%%%%%%%%%%%%%%%%%%%%%%%%%%%%%%%%%%%%%%%%%%%%%%%%%%%%%%%%%%
\subsection{Constitutive Modeling Principles}
\label{subsec:constitutive-modeling-principles}
%%%%%%%%%%%%%%%%%%%%%%%%%%%%%%%%%%%%%%%%%%%%%%%%%%%%%%%%%%%%

A constitutive law should respect:

\begin{itemize}
    \item material symmetry;
    \item frame indifference;
    \item dimensional consistency;
    \item thermodynamic restrictions;
    \item experimentally observed behavior;
    \item appropriate deformation scales.
\end{itemize}

A mathematically convenient stress-strain relation is not automatically
physically valid.

%%%%%%%%%%%%%%%%%%%%%%%%%%%%%%%%%%%%%%%%%%%%%%%%%%%%%%%%%%%%
\subsection{Material Symmetry}
\label{subsec:material-symmetry}
%%%%%%%%%%%%%%%%%%%%%%%%%%%%%%%%%%%%%%%%%%%%%%%%%%%%%%%%%%%%

An isotropic material has no preferred direction.

An anisotropic material has direction-dependent properties.

Examples include:

\[
\boxed{
\text{wood},
}
\]

\[
\boxed{
\text{fiber composites},
}
\]

and

\[
\boxed{
\text{biological tissues}.
}
\]

Material symmetry determines the allowed structure of constitutive tensors.

%%%%%%%%%%%%%%%%%%%%%%%%%%%%%%%%%%%%%%%%%%%%%%%%%%%%%%%%%%%%
\subsection{Elastic Anisotropy}
\label{subsec:elastic-anisotropy}
%%%%%%%%%%%%%%%%%%%%%%%%%%%%%%%%%%%%%%%%%%%%%%%%%%%%%%%%%%%%

In the most general linear elastic case,

\[
\boxed{
\sigma_{ij}
=
C_{ijkl}\varepsilon_{kl}.
}
\]

Symmetries of stress, strain, and elastic energy reduce the number of
independent coefficients.

Isotropy reduces the law to only two independent elastic constants.

%%%%%%%%%%%%%%%%%%%%%%%%%%%%%%%%%%%%%%%%%%%%%%%%%%%%%%%%%%%%
\subsection{Bulk Modulus}
\label{subsec:bulk-modulus}
%%%%%%%%%%%%%%%%%%%%%%%%%%%%%%%%%%%%%%%%%%%%%%%%%%%%%%%%%%%%

The bulk modulus measures resistance to volumetric deformation.

For small isotropic compression,

\[
\boxed{
K
=
-\frac{p}
{\Delta V/V}.
}
\]

For isotropic linear elasticity,

\[
\boxed{
K
=
\lambda+\frac{2}{3}\mu.
}
\]

Large \(K\) corresponds to low compressibility.

%%%%%%%%%%%%%%%%%%%%%%%%%%%%%%%%%%%%%%%%%%%%%%%%%%%%%%%%%%%%
\subsection{Volumetric Strain}
\label{subsec:volumetric-strain}
%%%%%%%%%%%%%%%%%%%%%%%%%%%%%%%%%%%%%%%%%%%%%%%%%%%%%%%%%%%%

For small deformation,

\[
\boxed{
\varepsilon_v
=
\operatorname{tr}
\boldsymbol{\varepsilon}.
}
\]

This approximates the fractional volume change:

\[
\boxed{
\varepsilon_v
\approx
\frac{\Delta V}{V}.
}
\]

For incompressible small-strain deformation,

\[
\boxed{
\operatorname{tr}
\boldsymbol{\varepsilon}=0.
}
\]

%%%%%%%%%%%%%%%%%%%%%%%%%%%%%%%%%%%%%%%%%%%%%%%%%%%%%%%%%%%%
\subsection{Deviatoric Strain}
\label{subsec:deviatoric-strain}
%%%%%%%%%%%%%%%%%%%%%%%%%%%%%%%%%%%%%%%%%%%%%%%%%%%%%%%%%%%%

The strain tensor may be decomposed as

\[
\boxed{
\boldsymbol{\varepsilon}
=
\frac13
\operatorname{tr}
(\boldsymbol{\varepsilon})
\mathbf{I}
+
\mathbf{e},
}
\]

where

\[
\boxed{
\operatorname{tr}\mathbf{e}=0.
}
\]

The tensor \(\mathbf{e}\) represents distortional deformation without volume
change.

%%%%%%%%%%%%%%%%%%%%%%%%%%%%%%%%%%%%%%%%%%%%%%%%%%%%%%%%%%%%
\subsection{One-Dimensional Wave Equation in a Bar}
\label{subsec:bar-wave-equation}
%%%%%%%%%%%%%%%%%%%%%%%%%%%%%%%%%%%%%%%%%%%%%%%%%%%%%%%%%%%%

For a slender elastic bar with displacement \(u(x,t)\),

\[
\boxed{
\rho A u_{tt}
=
\frac{\partial}{\partial x}
\left(
EAu_x
\right).
}
\]

For constant \(E\), \(A\), and \(\rho\),

\[
\boxed{
u_{tt}
=
c^2u_{xx},
}
\]

where

\[
\boxed{
c
=
\sqrt{\frac{E}{\rho}}.
}
\]

Thus elastic disturbances propagate as waves.

%%%%%%%%%%%%%%%%%%%%%%%%%%%%%%%%%%%%%%%%%%%%%%%%%%%%%%%%%%%%
\subsection{Worked Example: Bar Wave Speed}
\label{subsec:worked-bar-wave-speed}
%%%%%%%%%%%%%%%%%%%%%%%%%%%%%%%%%%%%%%%%%%%%%%%%%%%%%%%%%%%%

\begin{workedexamplebox}{Wave Speed in a Uniform Elastic Bar}

Suppose

\[
E=200\times10^9\text{ Pa}
\]

and

\[
\rho=7800\text{ kg/m}^3.
\]

Then

\[
c
=
\sqrt{
\frac{200\times10^9}
{7800}
}.
\]

Numerically,

\[
c
\approx
5064\text{ m/s}.
\]

\begin{answerbox}
\[
\boxed{
c\approx5.06\times10^3\text{ m/s}.
}
\]
\end{answerbox}

\end{workedexamplebox}

%%%%%%%%%%%%%%%%%%%%%%%%%%%%%%%%%%%%%%%%%%%%%%%%%%%%%%%%%%%%
\subsection{Continuum Mechanics Workflow}
\label{subsec:continuum-mechanics-workflow}
%%%%%%%%%%%%%%%%%%%%%%%%%%%%%%%%%%%%%%%%%%%%%%%%%%%%%%%%%%%%

A useful modeling workflow is:

\begin{enumerate}
    \item identify the body and reference configuration;
    \item define the motion map;
    \item determine whether a material or spatial description is more useful;
    \item define displacement, velocity, and acceleration;
    \item compute the deformation gradient;
    \item compute the Jacobian;
    \item choose an appropriate strain measure;
    \item identify material density;
    \item impose mass conservation;
    \item identify body and surface forces;
    \item define the stress measure;
    \item apply balance of linear momentum;
    \item enforce angular momentum balance;
    \item choose a constitutive law;
    \item verify material symmetry;
    \item verify dimensional consistency;
    \item determine whether small-strain assumptions are appropriate;
    \item determine whether incompressibility is appropriate;
    \item include thermal effects when needed;
    \item impose boundary conditions;
    \item impose initial conditions for dynamic problems;
    \item nondimensionalize when useful;
    \item derive weak forms when numerical approximation is required;
    \item solve analytically or numerically;
    \item verify conservation and boundary conditions;
    \item compare predicted deformation and stress with experiments;
    \item revise the constitutive model if necessary.
\end{enumerate}

%%%%%%%%%%%%%%%%%%%%%%%%%%%%%%%%%%%%%%%%%%%%%%%%%%%%%%%%%%%%
\subsection{Common Errors}
\label{subsec:continuum-mechanics-common-errors}
%%%%%%%%%%%%%%%%%%%%%%%%%%%%%%%%%%%%%%%%%%%%%%%%%%%%%%%%%%%%

\subsubsection{Treating Matter as Literally Continuous at All Scales}

The continuum hypothesis is an approximation.

At sufficiently small scales, molecular or atomistic models may be required.

%%%%%%%%%%%%%%%%%%%%%%%%%%%%%%%%%%%%%%%%%%%%%%%%%%%%%%%%%%%%
\subsubsection{Confusing Reference and Current Coordinates}

\[
\mathbf{X}
\]

labels a material point in the reference configuration.

\[
\mathbf{x}
\]

is its current spatial location.

They should not be used interchangeably.

%%%%%%%%%%%%%%%%%%%%%%%%%%%%%%%%%%%%%%%%%%%%%%%%%%%%%%%%%%%%
\subsubsection{Confusing Displacement With Position}

Displacement is

\[
\boxed{
\mathbf{u}
=
\mathbf{x}-\mathbf{X}.
}
\]

It measures change in position, not the position itself.

%%%%%%%%%%%%%%%%%%%%%%%%%%%%%%%%%%%%%%%%%%%%%%%%%%%%%%%%%%%%
\subsubsection{Confusing Eulerian and Lagrangian Derivatives}

In the Eulerian description, a moving particle experiences the material
derivative

\[
\frac{D}{Dt}
=
\frac{\partial}{\partial t}
+
\mathbf{v}\cdot\nabla.
\]

The ordinary partial derivative at fixed spatial location is generally not
the same.

%%%%%%%%%%%%%%%%%%%%%%%%%%%%%%%%%%%%%%%%%%%%%%%%%%%%%%%%%%%%
\subsubsection{Forgetting the Convective Acceleration}

The acceleration field is

\[
\boxed{
\mathbf{a}
=
\frac{\partial\mathbf{v}}{\partial t}
+
(\mathbf{v}\cdot\nabla)\mathbf{v}.
}
\]

Even a steady velocity field can have nonzero acceleration.

%%%%%%%%%%%%%%%%%%%%%%%%%%%%%%%%%%%%%%%%%%%%%%%%%%%%%%%%%%%%
\subsubsection{Confusing Deformation Gradient With Strain}

The deformation gradient

\[
\mathbf{F}
\]

contains both deformation and rotation.

A strain tensor is constructed to isolate deformation measures.

%%%%%%%%%%%%%%%%%%%%%%%%%%%%%%%%%%%%%%%%%%%%%%%%%%%%%%%%%%%%
\subsubsection{Assuming \(\mathbf{F}\) Is Symmetric}

The deformation gradient need not be symmetric.

It may contain rotational components.

%%%%%%%%%%%%%%%%%%%%%%%%%%%%%%%%%%%%%%%%%%%%%%%%%%%%%%%%%%%%
\subsubsection{Using Infinitesimal Strain for Large Deformation}

The approximation

\[
\boldsymbol{\varepsilon}
=
\frac12
\left(
\nabla\mathbf{u}
+
\nabla\mathbf{u}^{T}
\right)
\]

is valid only for sufficiently small displacement gradients.

%%%%%%%%%%%%%%%%%%%%%%%%%%%%%%%%%%%%%%%%%%%%%%%%%%%%%%%%%%%%
\subsubsection{Confusing Stress and Strain}

Stress has units of force per area.

Strain is dimensionless.

Stress describes internal forces.

Strain describes deformation.

%%%%%%%%%%%%%%%%%%%%%%%%%%%%%%%%%%%%%%%%%%%%%%%%%%%%%%%%%%%%
\subsubsection{Confusing Pressure With General Stress}

Pressure corresponds to the isotropic part of stress.

A general solid stress tensor may also contain shear and deviatoric
components.

%%%%%%%%%%%%%%%%%%%%%%%%%%%%%%%%%%%%%%%%%%%%%%%%%%%%%%%%%%%%
\subsubsection{Assuming the Stress Tensor Is Always Diagonal}

Stress can contain shear components.

It becomes diagonal only in principal directions or under special loading.

%%%%%%%%%%%%%%%%%%%%%%%%%%%%%%%%%%%%%%%%%%%%%%%%%%%%%%%%%%%%
\subsubsection{Ignoring Stress Symmetry Conditions}

For an ordinary Cauchy continuum without couple stresses,

\[
\boldsymbol{\sigma}
=
\boldsymbol{\sigma}^{T}.
\]

More generalized continuum theories may modify this conclusion.

%%%%%%%%%%%%%%%%%%%%%%%%%%%%%%%%%%%%%%%%%%%%%%%%%%%%%%%%%%%%
\subsubsection{Confusing Body Forces and Surface Forces}

Body forces scale with volume.

Surface tractions scale with area.

They enter momentum balance differently.

%%%%%%%%%%%%%%%%%%%%%%%%%%%%%%%%%%%%%%%%%%%%%%%%%%%%%%%%%%%%
\subsubsection{Using Hooke's Law Beyond Its Valid Range}

Linear elasticity is an approximation.

Many materials become nonlinear, plastic, damaged, or fractured at large
loads.

%%%%%%%%%%%%%%%%%%%%%%%%%%%%%%%%%%%%%%%%%%%%%%%%%%%%%%%%%%%%
\subsubsection{Assuming Isotropy}

Many materials have direction-dependent properties.

An isotropic constitutive law may be inappropriate for composites, crystals,
wood, or biological tissue.

%%%%%%%%%%%%%%%%%%%%%%%%%%%%%%%%%%%%%%%%%%%%%%%%%%%%%%%%%%%%
\subsubsection{Confusing Elasticity and Plasticity}

Elastic deformation is recoverable.

Plastic deformation is irreversible.

A material may exhibit both depending on stress level.

%%%%%%%%%%%%%%%%%%%%%%%%%%%%%%%%%%%%%%%%%%%%%%%%%%%%%%%%%%%%
\subsubsection{Confusing Elasticity and Viscoelasticity}

Elastic stress depends on deformation.

Viscoelastic response also depends on time or deformation rate.

%%%%%%%%%%%%%%%%%%%%%%%%%%%%%%%%%%%%%%%%%%%%%%%%%%%%%%%%%%%%
\subsubsection{Ignoring Thermal Strain}

Temperature changes can generate significant deformation and stress,
especially under constrained expansion.

%%%%%%%%%%%%%%%%%%%%%%%%%%%%%%%%%%%%%%%%%%%%%%%%%%%%%%%%%%%%
\subsubsection{Ignoring Boundary Conditions}

The governing PDEs alone usually do not determine a unique continuum
solution.

Boundary conditions are essential.

%%%%%%%%%%%%%%%%%%%%%%%%%%%%%%%%%%%%%%%%%%%%%%%%%%%%%%%%%%%%
\subsubsection{Ignoring Initial Conditions in Dynamic Problems}

Elastic wave and transient mechanics problems require initial displacement
and velocity.

%%%%%%%%%%%%%%%%%%%%%%%%%%%%%%%%%%%%%%%%%%%%%%%%%%%%%%%%%%%%
\subsubsection{Confusing Material Incompressibility With Constant Density in Every Context}

For an incompressible material,

\[
J=1.
\]

In an Eulerian constant-density flow, this leads to

\[
\nabla\cdot\mathbf{v}=0.
\]

The exact interpretation depends on the model and variables used.

%%%%%%%%%%%%%%%%%%%%%%%%%%%%%%%%%%%%%%%%%%%%%%%%%%%%%%%%%%%%
\subsubsection{Assuming Numerical Convergence Guarantees Physical Accuracy}

A finite element solution may converge accurately to the equations while the
constitutive law itself poorly represents the actual material.

%%%%%%%%%%%%%%%%%%%%%%%%%%%%%%%%%%%%%%%%%%%%%%%%%%%%%%%%%%%%
\subsection{Applications}
\label{subsec:continuum-mechanics-applications}
%%%%%%%%%%%%%%%%%%%%%%%%%%%%%%%%%%%%%%%%%%%%%%%%%%%%%%%%%%%%

\begin{applicationbox}

\textbf{Structural Mechanics.}
Stress, strain, elasticity, and finite element methods are used to analyze
buildings, bridges, aircraft, and mechanical components.

\medskip

\textbf{Solid Mechanics.}
Elasticity, plasticity, fracture, and wave propagation describe deformation
of solid materials.

\medskip

\textbf{Fluid Mechanics.}
The same conservation laws of mass, momentum, and energy form the basis of
fluid dynamics.

\medskip

\textbf{Geophysics.}
Elastic and viscoelastic continuum models describe seismic waves,
deformation of Earth's crust, and mantle processes.

\medskip

\textbf{Biomechanics.}
Soft tissues, bones, blood vessels, and muscles are modeled as deformable
continua.

\medskip

\textbf{Materials Science.}
Constitutive models connect microstructure with macroscopic stress,
deformation, yielding, and thermal response.

\medskip

\textbf{Acoustics and Elastic Waves.}
Continuum equations describe longitudinal and transverse wave propagation.

\medskip

\textbf{Computational Mechanics.}
Finite element, finite difference, and spectral methods approximate complex
continuum systems computationally.

\end{applicationbox}

%%%%%%%%%%%%%%%%%%%%%%%%%%%%%%%%%%%%%%%%%%%%%%%%%%%%%%%%%%%%
\subsection{Exercises}
\label{subsec:continuum-mechanics-exercises}
%%%%%%%%%%%%%%%%%%%%%%%%%%%%%%%%%%%%%%%%%%%%%%%%%%%%%%%%%%%%

\begin{exercise}[1]
Explain the continuum hypothesis.
\end{exercise}

\begin{exercise}[1]
Distinguish reference and current configurations.
\end{exercise}

\begin{exercise}[1]
Define the motion map

\[
\mathbf{x}
=
\boldsymbol{\chi}(\mathbf{X},t).
\]
\end{exercise}

\begin{exercise}[1]
Define displacement.
\end{exercise}

\begin{exercise}[1]
Define the material derivative.
\end{exercise}

\begin{exercise}[1]
Define the deformation gradient.
\end{exercise}

\begin{exercise}[1]
Define the Jacobian

\[
J=\det\mathbf{F}.
\]
\end{exercise}

\begin{exercise}[1]
Define infinitesimal strain.
\end{exercise}

\begin{exercise}[1]
Define the Cauchy stress tensor.
\end{exercise}

\begin{exercise}[1]
State Cauchy's stress theorem.
\end{exercise}

\begin{exercise}[2]
For

\[
x_1=2X_1,
\qquad
x_2=3X_2,
\]

compute the deformation gradient.
\end{exercise}

\begin{exercise}[2]
For the same deformation, compute \(J\).
\end{exercise}

\begin{exercise}[2]
Interpret the value of \(J\) geometrically.
\end{exercise}

\begin{exercise}[2]
For displacement

\[
\mathbf{u}(x,y)
=
\begin{bmatrix}
ax\\
by
\end{bmatrix},
\]

compute the infinitesimal strain tensor.
\end{exercise}

\begin{exercise}[2]
For velocity field

\[
\mathbf{v}
=
\begin{bmatrix}
ax\\
-by
\end{bmatrix},
\]

compute

\[
\nabla\cdot\mathbf{v}.
\]
\end{exercise}

\begin{exercise}[3]
Determine when the preceding flow is incompressible.
\end{exercise}

\begin{exercise}[2]
Compute the material derivative of

\[
f(x,y,t)=x^2+y+t
\]

for velocity field

\[
\mathbf{v}=(1,2).
\]
\end{exercise}

\begin{exercise}[2]
Show that rigid translation produces zero strain.
\end{exercise}

\begin{exercise}[3]
Show that finite rigid rotation gives

\[
\mathbf{C}=\mathbf{I}.
\]
\end{exercise}

\begin{exercise}[3]
Explain why linearized strain can incorrectly interpret large rigid
rotations.
\end{exercise}

\begin{exercise}[2]
Write the continuity equation.
\end{exercise}

\begin{exercise}[3]
Derive

\[
\frac{D\rho}{Dt}
+
\rho\nabla\cdot\mathbf{v}
=
0
\]

from

\[
\rho_t+\nabla\cdot(\rho\mathbf{v})=0.
\]
\end{exercise}

\begin{exercise}[3]
Show that

\[
\rho J=\rho_0
\]

expresses mass conservation.
\end{exercise}

\begin{exercise}[2]
Define a traction vector.
\end{exercise}

\begin{exercise}[2]
Given

\[
\boldsymbol{\sigma}
=
\begin{bmatrix}
4 & 1\\
1 & 2
\end{bmatrix}
\]

and

\[
\mathbf{n}
=
\begin{bmatrix}
1\\
0
\end{bmatrix},
\]

compute the traction.
\end{exercise}

\begin{exercise}[3]
For the same stress tensor, compute the principal stresses.
\end{exercise}

\begin{exercise}[3]
Find the associated principal directions.
\end{exercise}

\begin{exercise}[2]
State the balance of linear momentum.
\end{exercise}

\begin{exercise}[2]
State the equilibrium equation for a static continuum.
\end{exercise}

\begin{exercise}[3]
Explain why angular momentum balance implies stress symmetry in a classical
continuum.
\end{exercise}

\begin{exercise}[2]
Decompose

\[
\boldsymbol{\sigma}
\]

into hydrostatic and deviatoric parts.
\end{exercise}

\begin{exercise}[3]
Show that the deviatoric stress has zero trace.
\end{exercise}

\begin{exercise}[2]
State the isotropic linear elastic constitutive law.
\end{exercise}

\begin{exercise}[2]
Explain the physical meaning of Young's modulus.
\end{exercise}

\begin{exercise}[2]
Explain Poisson ratio.
\end{exercise}

\begin{exercise}[3]
Convert between \((E,\nu)\) and \((\lambda,\mu)\).
\end{exercise}

\begin{exercise}[3]
Derive the one-dimensional elastic extension

\[
\Delta L
=
\frac{\sigma L}{E}.
\]
\end{exercise}

\begin{exercise}[3]
Derive the Navier--Cauchy equations from momentum balance and isotropic
Hooke's law.
\end{exercise}

\begin{exercise}[2]
Define elastic strain-energy density.
\end{exercise}

\begin{exercise}[3]
Show that for one-dimensional linear elasticity,

\[
W
=
\frac12E\varepsilon^2.
\]
\end{exercise}

\begin{exercise}[2]
Define hyperelasticity conceptually.
\end{exercise}

\begin{exercise}[2]
Define frame indifference conceptually.
\end{exercise}

\begin{exercise}[2]
Write the Kelvin--Voigt constitutive law.
\end{exercise}

\begin{exercise}[2]
Write the Maxwell constitutive law.
\end{exercise}

\begin{exercise}[3]
Solve the Kelvin--Voigt creep problem for constant stress.
\end{exercise}

\begin{exercise}[3]
Study stress relaxation in the Maxwell model.
\end{exercise}

\begin{exercise}[2]
Define plastic deformation.
\end{exercise}

\begin{exercise}[2]
Define a yield surface conceptually.
\end{exercise}

\begin{exercise}[3]
Explain the decomposition

\[
\boldsymbol{\varepsilon}
=
\boldsymbol{\varepsilon}^{e}
+
\boldsymbol{\varepsilon}^{p}.
\]
\end{exercise}

\begin{exercise}[2]
Compute thermal strain for a material with coefficient \(\alpha\) and
temperature increase \(\Delta T\).
\end{exercise}

\begin{exercise}[3]
Explain why constrained thermal expansion generates stress.
\end{exercise}

\begin{exercise}[2]
Write Fourier's law of heat conduction.
\end{exercise}

\begin{exercise}[3]
Explain the mechanical meaning of

\[
\boldsymbol{\sigma}:\mathbf{D}.
\]
\end{exercise}

\begin{exercise}[3]
Explain why thermodynamic restrictions are required on constitutive laws.
\end{exercise}

\begin{exercise}[2]
Distinguish displacement and traction boundary conditions.
\end{exercise}

\begin{exercise}[3]
Explain the principle of virtual work.
\end{exercise}

\begin{exercise}[3]
Explain why weak formulations are useful for finite element methods.
\end{exercise}

\begin{exercise}[2]
State the longitudinal elastic wave speed.
\end{exercise}

\begin{exercise}[2]
State the shear-wave speed.
\end{exercise}

\begin{exercise}[3]
Show that

\[
c_P>c_S
\]

for ordinary positive Lamé parameters.
\end{exercise}

\begin{exercise}[3]
Derive the one-dimensional elastic bar wave equation.
\end{exercise}

\begin{exercise}[2]
Determine the dimensions of stress.
\end{exercise}

\begin{exercise}[2]
Show that strain is dimensionless.
\end{exercise}

\begin{exercise}[3]
Explain why Young's modulus has the same dimensions as stress.
\end{exercise}

\begin{exercise}[3]
Explain the distinction between model error and discretization error in a
finite element simulation.
\end{exercise}

\begin{exercise}[4]
Derive the finite-deformation mass relation

\[
\rho J=\rho_0.
\]
\end{exercise}

\begin{exercise}[4]
Study the polar decomposition

\[
\mathbf{F}
=
\mathbf{R}\mathbf{U}
=
\mathbf{V}\mathbf{R}.
\]
\end{exercise}

\begin{exercise}[4]
Relate the stretch tensors to the Cauchy--Green tensors.
\end{exercise}

\begin{exercise}[4]
Study principal stretches.
\end{exercise}

\begin{exercise}[4]
Derive Cauchy's stress theorem from force balance on an infinitesimal
tetrahedron.
\end{exercise}

\begin{exercise}[4]
Derive stress symmetry from angular momentum balance.
\end{exercise}

\begin{exercise}[4]
Study Mohr's circle for a two-dimensional stress state.
\end{exercise}

\begin{exercise}[4]
Study isotropic elasticity in index notation.
\end{exercise}

\begin{exercise}[4]
Derive elastic-wave equations using Helmholtz decomposition.
\end{exercise}

\begin{exercise}[4]
Study hyperelastic constitutive models.
\end{exercise}

\begin{exercise}[4]
Study incompressible hyperelasticity.
\end{exercise}

\begin{exercise}[4]
Study finite-strain stress measures.
\end{exercise}

\begin{exercise}[4]
Study von Mises yielding.
\end{exercise}

\begin{exercise}[4]
Study Tresca yielding.
\end{exercise}

\begin{exercise}[4]
Study linear viscoelastic hereditary integrals.
\end{exercise}

\begin{exercise}[4]
Derive a thermoelastic constitutive model.
\end{exercise}

\begin{exercise}[4]
Derive the weak form of linear elasticity.
\end{exercise}

\begin{exercise}[4]
Construct the stiffness matrix of a one-dimensional finite element.
\end{exercise}

\begin{exercise}[4]
Develop a numerical solution for a loaded elastic bar.
\end{exercise}

%%%%%%%%%%%%%%%%%%%%%%%%%%%%%%%%%%%%%%%%%%%%%%%%%%%%%%%%%%%%
\subsection{Conceptual Summary}
\label{subsec:continuum-mechanics-summary}
%%%%%%%%%%%%%%%%%%%%%%%%%%%%%%%%%%%%%%%%%%%%%%%%%%%%%%%%%%%%

Continuum mechanics models matter as continuously distributed through space.

A material point moves according to

\[
\boxed{
\mathbf{x}
=
\boldsymbol{\chi}(\mathbf{X},t).
}
\]

Local deformation is described by

\[
\boxed{
\mathbf{F}
=
\frac{\partial\mathbf{x}}
{\partial\mathbf{X}},
}
\]

with volume ratio

\[
\boxed{
J=\det\mathbf{F}.
}
\]

For small deformation, strain is approximated by

\[
\boxed{
\boldsymbol{\varepsilon}
=
\frac12
\left(
\nabla\mathbf{u}
+
(\nabla\mathbf{u})^T
\right).
}
\]

Mass conservation gives

\[
\boxed{
\rho_t
+
\nabla\cdot(\rho\mathbf{v})
=
0.
}
\]

Linear momentum balance gives

\[
\boxed{
\rho
\frac{D\mathbf{v}}{Dt}
=
\nabla\cdot\boldsymbol{\sigma}
+
\rho\mathbf{b}.
}
\]

Surface traction satisfies

\[
\boxed{
\mathbf{t}
=
\boldsymbol{\sigma}\mathbf{n}.
}
\]

Material behavior is supplied by constitutive laws.

For isotropic linear elasticity,

\[
\boxed{
\boldsymbol{\sigma}
=
\lambda
\operatorname{tr}
(\boldsymbol{\varepsilon})
\mathbf{I}
+
2\mu\boldsymbol{\varepsilon}.
}
\]

Continuum mechanics therefore combines

\[
\boxed{
\text{kinematics}
+
\text{conservation laws}
+
\text{constitutive behavior}
}
\]

to produce field equations for deformable matter.

These equations form the mathematical foundation of solid mechanics,
elasticity, plasticity, viscoelasticity, biomechanics, geophysics, and fluid
mechanics.

%%%%%%%%%%%%%%%%%%%%%%%%%%%%%%%%%%%%%%%%%%%%%%%%%%%%%%%%%%%%
\subsection{Section Connections}
\label{subsec:continuum-mechanics-connections}
%%%%%%%%%%%%%%%%%%%%%%%%%%%%%%%%%%%%%%%%%%%%%%%%%%%%%%%%%%%%

\chapterconnection{
Continuum Mechanics extends the field-based ideas of Mathematical Physics
from particles and abstract fields to continuously distributed matter.
Vector calculus describes local transport and force balance, tensor methods
describe stress and deformation, partial differential equations govern
spatially distributed motion, and constitutive laws connect mathematical
kinematics with physical material behavior. Linear elasticity provides the
basic solid-mechanics example, while finite deformation, viscoelasticity,
plasticity, thermal effects, and numerical weak formulations extend the
framework. The next section, Fluid Dynamics, specializes continuum mechanics
to fluids, where stress is strongly connected to pressure and rate of
deformation and the conservation laws lead to the Euler and Navier--Stokes
equations.
}

%%%%%%%%%%%%%%%%%%%%%%%%%%%%%%%%%%%%%%%%%%%%%%%%%%%%%%%%%%%%
% SECTION SOURCE TRACKING
%%%%%%%%%%%%%%%%%%%%%%%%%%%%%%%%%%%%%%%%%%%%%%%%%%%%%%%%%%%%

%------------------------------------------------------------
% 13.4 Continuum Mechanics
%------------------------------------------------------------
%
% Source basis:
%   - Standard continuum mechanics.
%   - Standard elasticity and solid mechanics.
%   - Standard finite-deformation kinematics.
%   - Standard balance laws and constitutive modeling.
%
% Used for:
%   - continuum hypothesis;
%   - material points;
%   - reference and current configurations;
%   - motion maps;
%   - displacement;
%   - Lagrangian and Eulerian descriptions;
%   - velocity and acceleration;
%   - material derivatives;
%   - deformation gradients;
%   - Jacobians;
%   - incompressibility;
%   - Cauchy--Green tensors;
%   - Green--Lagrange strain;
%   - infinitesimal strain;
%   - normal and shear strain;
%   - rigid-body motions;
%   - rate-of-deformation and spin tensors;
%   - mass conservation;
%   - density transformation;
%   - body and surface forces;
%   - traction;
%   - Cauchy stress;
%   - momentum balance;
%   - angular momentum balance;
%   - stress symmetry;
%   - principal stress;
%   - hydrostatic and deviatoric stress;
%   - constitutive equations;
%   - linear elasticity;
%   - isotropic elasticity;
%   - Lamé parameters;
%   - Young's modulus;
%   - Poisson ratio;
%   - Navier--Cauchy equations;
%   - elastic strain energy;
%   - hyperelasticity;
%   - objectivity;
%   - viscoelasticity;
%   - Kelvin--Voigt models;
%   - Maxwell models;
%   - plasticity;
%   - yield criteria;
%   - thermal strain;
%   - thermoelasticity;
%   - energy balance;
%   - heat conduction;
%   - thermodynamic restrictions;
%   - boundary and initial conditions;
%   - elastic waves;
%   - longitudinal and transverse waves;
%   - weak formulations;
%   - virtual work;
%   - finite element preview;
%   - stiffness matrices;
%   - dimensional analysis;
%   - small-strain approximation;
%   - finite-deformation stress measures;
%   - material symmetry;
%   - anisotropy;
%   - bulk modulus;
%   - volumetric and deviatoric strain;
%   - bar-wave equations;
%   - applications and exercises.
%
% Notes:
%   - Fluid Dynamics is developed next in Section 13.5.
%   - Control Theory follows in Section 13.6.
%   - Network Science follows in Section 13.7.
%   - Mathematical Finance follows in Section 13.8.
%   - Mathematical Economics follows in Section 13.9.
%
%%%%%%%%%%%%%%%%%%%%%%%%%%%%%%%%%%%%%%%%%%%%%%%%%%%%%%%%%%%%